\chapter{IFRS Recognition and Measurement Mapping}
\label{ch:ifrs}

This chapter provides paragraph-level mapping between IFRS standards and AccDSL
recognition, measurement, and presentation operations. It addresses the gap
identified in v1.2 where only "orientation overviews" were used. Each mapping
includes the standard paragraph reference, economic preconditions, accounting
operation, state changes, and conformance tests.

\section{Mapping methodology and authority}

Every AccDSL accounting operation claimed to implement IFRS must provide:

\begin{enumerate}
\item \textbf{Standard reference:} IFRS/IAS number, paragraph(s), effective date,
and any relevant amendments or interpretations (IFRIC/SIC).
\item \textbf{Scope determination:} Which entity types, instruments, or transactions
the operation applies to, with explicit out-of-scope cases.
\item \textbf{Economic preconditions:} Contract terms, market conditions, or
entity characteristics that must be verified before the rule applies.
\item \textbf{Recognition criteria:} When an asset, liability, income, or expense
is recognized, with timing and measurement basis.
\item \textbf{Subsequent measurement:} Amortized cost, fair value through profit
or loss (FVTPL), fair value through OCI (FVOCI), or other basis, with
remeasurement triggers.
\item \textbf{Presentation and disclosure:} Required line items, note schedules,
and movement reconciliations.
\item \textbf{Implementation contract:} AccDSL event types, rule signatures,
state fields, and journal templates.
\item \textbf{Conformance tests:} Worked examples with known inputs and expected
outputs, including edge cases.
\end{enumerate}

\subsection{Authority and review status}

The mappings in this chapter are based on:
\begin{itemize}
\item IFRS Standards consolidated text (Red Book 2024 edition) where accessible
\item IFRS Foundation website official summaries and illustrative examples
\item AASB texts incorporating IAS/IFRS with Australian-specific paragraphs noted
\item Big Four firm technical guidance (PwC, Deloitte, EY, KPMG) for worked examples
\end{itemize}

\warn{Review status as of 5 October 2026:} This chapter provides design-level
mapping. \textbf{Independent CPA/accounting-firm review is required} before any
operation can be certified as IFRS-compliant. Current status is marked for each
operation: \done{Reviewed}, \warn{Pending Review}, or \warn{Not Implemented}.

\section{Financial instruments: IFRS 9}

\subsection{Scope and classification (IFRS 9.2.1--9.4.3)}

IFRS 9 applies to financial assets and financial liabilities, with specific scope
exclusions for leases (IFRS 16), insurance contracts (IFRS 17), employee benefits
(IAS 19), and equity instruments issued by the entity (IAS 32).

\textbf{Classification decision tree:}

\begin{enumerate}
\item \textbf{Business model test (9.4.1):} How the entity manages financial assets:
  \begin{itemize}
  \item Hold to collect contractual cash flows $\Rightarrow$ amortized cost (if SPPI)
  \item Hold to collect and sell $\Rightarrow$ FVOCI (if SPPI)
  \item Other/trading $\Rightarrow$ FVTPL
  \end{itemize}
\item \textbf{Cash flow characteristics test (9.4.2, B4.1):} Solely payments of
  principal and interest (SPPI):
  \begin{itemize}
  \item Plain vanilla loan/bond: SPPI satisfied
  \item Equity-linked, commodity-linked, leveraged: SPPI fails $\Rightarrow$ FVTPL
  \end{itemize}
\item \textbf{Fair value option (9.4.4):} Entity may irrevocably designate as FVTPL
  to eliminate accounting mismatch.
\end{enumerate}

\textbf{AccDSL implementation contract:}

\begin{lstlisting}
event ClassifyFinancialAsset(
  instrument: Instrument,
  business_model: Enum[HoldToCollect, HoldAndSell, Other],
  sppi_test: Boolean,
  fv_option_elected: Boolean
) book EntityBook;

rule ifrs9_classification for ClassifyFinancialAsset as e {
  let category = match (e.business_model, e.sppi_test, e.fv_option_elected) {
    (HoldToCollect, true, false) => AmortizedCost,
    (HoldAndSell, true, false) => FVOCI,
    (_, _, true) => FVTPL,
    _ => FVTPL
  };
  update_instrument_measurement_category(e.instrument, category);
  post { /* No immediate posting; affects subsequent measurement */ }
}
\end{lstlisting}

\warn{Status: Pending Review.} Business-model assessment requires evidence of
past sales, management intent, and performance evaluation. SPPI test for complex
instruments (stepped interest, early-redemption features) requires detailed
contractual analysis not fully implemented.

\subsection{Initial recognition and measurement (IFRS 9.5.1.1)}

Financial assets and liabilities are recognized when the entity becomes party to
the contractual provisions. Initial measurement is at fair value plus (for assets
not at FVTPL) transaction costs directly attributable to acquisition.

\textbf{AccDSL implementation:}

\begin{lstlisting}
event AcquireFinancialAsset(
  instrument: Instrument,
  principal: Money<C>,
  transaction_costs: Money<C>,
  fair_value: Money<C>,
  measurement_category: MeasurementCategory
) book EntityBook;

rule ifrs9_initial_recognition for AcquireFinancialAsset as e {
  require e.principal > 0;
  let initial_amount = match e.measurement_category {
    AmortizedCost | FVOCI => e.fair_value + e.transaction_costs,
    FVTPL => e.fair_value
  };
  post {
    debit FinancialAsset[e.instrument] initial_amount;
    credit Cash e.principal + e.transaction_costs;
  }
  // If e.principal != e.fair_value, recognize day-one gain/loss for FVTPL
  if e.measurement_category == FVTPL && e.principal != e.fair_value {
    let day_one_pnl = e.fair_value - e.principal;
    post {
      debit_or_credit PnL_FairValueGains day_one_pnl;
    }
  }
}
\end{lstlisting}

\done{Status: Core logic reviewed.} Transaction-cost allocation for portfolios
and embedded-derivative separation (9.4.3) require additional operations.

\subsection{Effective interest method (IFRS 9.5.4.1, B5.4.1--B5.4.7)}

For assets/liabilities at amortized cost or FVOCI, interest revenue/expense uses
the effective interest rate (EIR) that discounts future cash flows to the initial
carrying amount.

\textbf{Mathematics:} For initial amount $A_0$ and cash flows $CF_t$ at times $t_i$,
the EIR $r$ solves
\begin{equation}
A_0 = \sum_{i} \frac{CF_{t_i}}{(1+r)^{t_i}}.
\end{equation}
Interest revenue in period $[s,t)$ is $A_s \cdot r \cdot \Delta t$, where $\Delta t$
is the accrual fraction. The carrying amount updates as $A_t = A_s(1+r)^{\Delta t} - CF_{[s,t)}$.

\textbf{AccDSL implementation:}

\begin{lstlisting}
event AccrueInterest(
  instrument: Instrument,
  opening_carrying: Money<C>,
  eir: Rate<effective>,
  accrual_fraction: YearFraction
) book EntityBook;

rule ifrs9_effective_interest for AccrueInterest as e {
  require e.eir > 0 per year;
  let interest = e.opening_carrying * e.eir * e.accrual_fraction;
  let rounded_interest = roundMoney(interest, currencyScale(C), "half_even");
  post {
    debit FinancialAsset[e.instrument] rounded_interest;
    credit InterestRevenue rounded_interest;
  }
}
\end{lstlisting}

\done{Status: Reviewed.} Calculation correct for fixed cash flows. Variable-rate
instruments require reset events. Credit-impaired assets (purchased or subsequent)
use credit-adjusted EIR (9.5.4.1), which is a separate operation marked \warn{Not Implemented}.

\subsection{Expected credit loss model (IFRS 9.5.5.1--9.5.5.20)}

Impairment uses forward-looking expected credit losses (ECL). Three stages:

\begin{enumerate}
\item \textbf{Stage 1 (9.5.5.5):} 12-month ECL for assets without significant increase
in credit risk since initial recognition.
\item \textbf{Stage 2 (9.5.5.3):} Lifetime ECL for assets with significant credit
deterioration but not credit-impaired.
\item \textbf{Stage 3 (9.5.5.4):} Lifetime ECL for credit-impaired assets; interest
on net carrying (gross less allowance).
\end{enumerate}

\textbf{Significant increase in credit risk (9.5.5.9--9.5.5.11):}
\begin{itemize}
\item Change in default risk over expected life (not merely incurred loss)
\item Rebuttable presumption: $>30$ days past due
\item Qualitative and quantitative factors
\item Compare origination date and reporting date risk
\end{itemize}

\textbf{AccDSL implementation contract (simplified):}

\begin{lstlisting}
event MeasureECL(
  cohort: LoanCohort,
  stage: Enum[Stage1, Stage2, Stage3],
  gross_carrying: Money<C>,
  lifetime_pd: Real in [0,1],
  twelve_month_pd: Real in [0,1],
  lgd: Real in [0,1],
  discount_rate: Rate<effective>
) book EntityBook;

rule ifrs9_ecl for MeasureECL as e {
  let pd = if e.stage == Stage1 then e.twelve_month_pd else e.lifetime_pd;
  let ead = e.gross_carrying; // Simplified; should include undrawn commitments
  let expected_loss = ead * pd * e.lgd;
  let discounted_loss = expected_loss / (1 + e.discount_rate);

  let old_allowance = get_allowance(e.cohort);
  let allowance_change = discounted_loss - old_allowance;

  post {
    debit CreditLossExpense allowance_change;
    credit Allowance[e.cohort] allowance_change;
  }
}
\end{lstlisting}

\warn{Status: Not Implemented.} This is a simplified point estimate. Full implementation
requires:
\begin{itemize}
\item Probability-weighted scenarios (9.5.5.18)
\item Discounting to reporting date
\item Forward-looking information incorporation (9.5.5.17)
\item Staging transfers and their accounting effects
\item Modified financial assets (9.5.4.3)
\item Write-offs and subsequent recoveries (9.5.5.10)
\end{itemize}

\textbf{Conformance test needed:} IFRS 9 Implementation Guidance examples B5.5.40--B5.5.54,
covering corporate bonds, trade receivables, and revolving facilities.

\section{Financial instruments: IAS 32 presentation}

\subsection{Liability vs. equity classification (IAS 32.11--32.23)}

The issuer classifies an instrument as liability or equity based on substance,
not legal form. A financial liability requires delivering cash or another financial
asset, or exchanging under potentially unfavorable conditions. Equity is the
residual after deducting liabilities.

\textbf{Decision factors (IAS 32.16--32.20):}
\begin{itemize}
\item Contractual obligation to deliver cash: liability
\item Settleable in entity's own equity at fixed number of shares for fixed cash: equity
\item Settleable in variable number of shares to deliver fixed cash value: liability
\item Discretionary dividends with no obligation: equity
\item Contingent settlement provisions: assess whether within control
\end{itemize}

\textbf{Compound instruments (IAS 32.28--32.32, AG30--AG35):}

An instrument with both liability and equity components (e.g., convertible bond)
is split. The liability is measured first at present value of future cash flows;
equity is the residual issue proceeds.

\textbf{Example from Chapter 8 v1.2 (expanded):}

Issue a 2-year bond for 100, paying annual coupons of 5, redeemable at 100 or
convertible into 10 shares (holder's option). Comparable liability without
conversion yields 10\%.

\textbf{Calculation:}
\begin{align}
L_0 &= \frac{5}{1.1} + \frac{105}{1.1^2} = \frac{11050}{121} \approx 91.3223,\\
E_0 &= 100 - L_0 \approx 8.6777.
\end{align}

\textbf{AccDSL implementation:}

\begin{lstlisting}
event IssueCompoundInstrument(
  issue_proceeds: Money<C>,
  coupon_schedule: CashFlowSchedule,
  redemption_amount: Money<C>,
  conversion_ratio: Quantity<shares>,
  comparable_liability_rate: Rate<effective>
) book EntityBook;

rule ias32_compound_split for IssueCompoundInstrument as e {
  let liability_pv = present_value(
    e.coupon_schedule.append(e.redemption_amount),
    e.comparable_liability_rate
  );
  require liability_pv <= e.issue_proceeds;

  let equity_component = e.issue_proceeds - liability_pv;

  post {
    debit Cash e.issue_proceeds;
    credit BondsPayable liability_pv;
    credit EquityConversionOption equity_component;
  }

  record_conversion_terms(e.conversion_ratio, e.redemption_amount);
}
\end{lstlisting}

\done{Status: Reviewed.} Covers vanilla convertible. More complex terms (adjustable
conversion price, contingent conversion, embedded puts/calls) require additional
analysis (AG31--AG35) marked \warn{Pending Review}.

\subsection{Treasury shares (IAS 32.33--32.34)}

If entity reacquires own equity instruments (treasury shares), the consideration
paid reduces equity. No gain or loss in profit or loss.

\textbf{AccDSL implementation:}

\begin{lstlisting}
event RepurchaseOwnShares(
  quantity: Quantity<shares>,
  cash_paid: Money<C>
) book EntityBook;

rule ias32_treasury for RepurchaseOwnShares as e {
  require e.quantity > 0 shares;
  post {
    debit TreasuryStock e.cash_paid;
    credit Cash e.cash_paid;
  }
  update_shares_outstanding(subtract: e.quantity);
}
\end{lstlisting}

\done{Status: Reviewed.} Subsequent reissue (AG36) and cancellation (AG37) are
separate operations marked \warn{Not Implemented}.

\section{Revenue from contracts with customers: IFRS 15}

IFRS 15 establishes a five-step model for revenue recognition.

\subsection{The five steps (IFRS 15.9--15.72)}

\begin{enumerate}
\item \textbf{Identify the contract (15.9--15.21):} Approved, rights/obligations
  identifiable, commercial substance, collectability probable.
\item \textbf{Identify performance obligations (15.22--15.30):} Distinct goods/services
  or series of distinct goods/services.
\item \textbf{Determine transaction price (15.46--15.72):} Variable consideration,
  significant financing, non-cash, consideration payable to customer.
\item \textbf{Allocate price to performance obligations (15.73--15.90):} Relative
  stand-alone selling prices, with adjustments for discounts and variable consideration.
\item \textbf{Recognize revenue (15.31--15.45):} Over time (15.35--15.37) or at
  a point in time (15.38).
\end{enumerate}

\subsection{Over-time recognition (IFRS 15.35--15.37, 15.39--15.45)}

Revenue is recognized over time if one of:
\begin{itemize}
\item Customer simultaneously receives and consumes benefits (15.35(a))
\item Entity's performance creates/enhances asset customer controls (15.35(b))
\item Entity's performance creates asset with no alternative use and enforceable
  right to payment for performance to date (15.35(c))
\end{itemize}

Progress measurement (15.39--15.45): output methods (surveys, units delivered)
or input methods (costs incurred, labor hours).

\textbf{AccDSL implementation (construction contract example):}

\begin{lstlisting}
event RecognizeRevenueOverTime(
  contract: Contract,
  performance_obligation: PerformanceObligation,
  total_transaction_price: Money<C>,
  progress_to_date: Real in [0,1],
  costs_incurred_to_date: Money<C>,
  contract_asset_opening: Money<C>,
  billings_to_date: Money<C>
) book EntityBook;

rule ifrs15_over_time for RecognizeRevenueOverTime as e {
  require e.progress_to_date >= 0 && e.progress_to_date <= 1;

  let cumulative_revenue = e.total_transaction_price * e.progress_to_date;
  let previously_recognized = get_cumulative_revenue(e.performance_obligation);
  let current_period_revenue = cumulative_revenue - previously_recognized;

  // Contract asset = cumulative revenue - cumulative billings
  let contract_asset_closing = cumulative_revenue - e.billings_to_date;
  let contract_asset_change = contract_asset_closing - e.contract_asset_opening;

  post {
    debit ContractAsset contract_asset_change;
    credit Revenue current_period_revenue;
  }

  // Costs to fulfill (15.95--15.98) handled separately
  post {
    debit CostOfSales e.costs_incurred_to_date;
    credit ConstructionInProgress e.costs_incurred_to_date;
  }
}
\end{lstlisting}

\warn{Status: Pending Review.} Simplified input-method example. Full implementation
requires:
\begin{itemize}
\item Variable consideration constraints (15.56--15.58)
\item Contract modifications (15.18--15.21)
\item Onerous performance obligations (IAS 37)
\item Contract costs capitalization (15.95--15.104)
\item Principal vs. agent determination (15.B34--B38)
\item Right-of-return estimation (15.B20--B27)
\end{itemize}

\section{Leases: IFRS 16}

\subsection{Lease identification (IFRS 16.9--16.15)}

A contract contains a lease if it conveys the right to control use of an identified
asset for a period in exchange for consideration. Control requires:
\begin{itemize}
\item Right to obtain substantially all economic benefits (16.B21)
\item Right to direct use of the asset (16.B24--B30)
\end{itemize}

\subsection{Lessee accounting: initial recognition (IFRS 16.22--16.28)}

At commencement, recognize:
\begin{itemize}
\item \textbf{Right-of-use (ROU) asset:} Lease liability + initial direct costs +
  prepayments - incentives (16.24)
\item \textbf{Lease liability:} Present value of lease payments not yet paid,
  discounted at implicit rate or incremental borrowing rate (16.26)
\end{itemize}

\textbf{AccDSL implementation:}

\begin{lstlisting}
event CommenceLease(
  lease: Lease,
  payment_schedule: PaymentSchedule,
  discount_rate: Rate<effective>,
  initial_direct_costs: Money<C>,
  prepayments: Money<C>,
  lease_incentives: Money<C>
) book EntityBook;

rule ifrs16_lessee_initial for CommenceLease as e {
  let lease_liability = present_value(e.payment_schedule, e.discount_rate);
  let rou_asset = lease_liability + e.initial_direct_costs +
                  e.prepayments - e.lease_incentives;

  post {
    debit RightOfUseAsset[e.lease] rou_asset;
    credit LeaseLiability[e.lease] lease_liability;
    credit Cash e.initial_direct_costs + e.prepayments;
    debit Cash e.lease_incentives;
  }
}
\end{lstlisting}

\subsection{Subsequent measurement (IFRS 16.29--16.35)}

\textbf{Lease liability (16.36):}
\begin{itemize}
\item Increase by interest on liability (effective interest method)
\item Decrease by lease payments
\item Remeasure when payments change (16.39--16.43)
\end{itemize}

\textbf{ROU asset (16.30--16.35):}
\begin{itemize}
\item Cost model: depreciate over shorter of lease term and useful life
\item Impairment per IAS 36
\end{itemize}

\begin{lstlisting}
event AccrueLeaseInterest(
  lease: Lease,
  opening_liability: Money<C>,
  discount_rate: Rate<effective>,
  period: YearFraction
) book EntityBook;

rule ifrs16_interest for AccrueLeaseInterest as e {
  let interest = e.opening_liability * e.discount_rate * e.period;
  post {
    debit InterestExpense_Lease interest;
    credit LeaseLiability[e.lease] interest;
  }
}

event MakeLeasePayment(
  lease: Lease,
  payment: Money<C>
) book EntityBook;

rule ifrs16_payment for MakeLeasePayment as e {
  post {
    debit LeaseLiability[e.lease] e.payment;
    credit Cash e.payment;
  }
}

event DepreciateROU(
  lease: Lease,
  rou_asset_opening: Money<C>,
  depreciation_period: YearFraction,
  total_lease_term: YearFraction
) book EntityBook;

rule ifrs16_depreciation for DepreciateROU as e {
  let depreciation = (e.rou_asset_opening / e.total_lease_term) * e.depreciation_period;
  post {
    debit Depreciation_ROU depreciation;
    credit AccumulatedDepreciation_ROU[e.lease] depreciation;
  }
}
\end{lstlisting}

\done{Status: Core reviewed.} Variable lease payments (16.38), lease modifications
(16.44--16.46), and sale-leaseback (16.98--16.103) marked \warn{Not Implemented}.

\section{Employee benefits: IAS 19}

\subsection{Defined contribution plans (IAS 19.50--19.53)}

Entity's obligation limited to contributions. Recognize expense when services
rendered.

\begin{lstlisting}
event RecognizeDefinedContribution(
  contribution: Money<C>,
  period: Period
) book EntityBook;

rule ias19_dc for RecognizeDefinedContribution as e {
  post {
    debit EmployeeBenefitExpense e.contribution;
    credit PensionPayable e.contribution;
  }
}
\end{lstlisting}

\done{Status: Reviewed.}

\subsection{Defined benefit plans (IAS 19.54--19.152)}

Entity bears actuarial and investment risk. Present value of obligation (DBO)
and plan assets measured at each reporting date.

\textbf{Components (IAS 19.120--19.130):}
\begin{itemize}
\item \textbf{Service cost:} Current service, past service, gains/losses on settlement
  $\Rightarrow$ P\&L
\item \textbf{Net interest:} Net defined benefit liability/asset $\times$ discount rate
  $\Rightarrow$ P\&L
\item \textbf{Remeasurements:} Actuarial gains/losses, return on assets (excl. interest),
  asset ceiling changes $\Rightarrow$ OCI
\end{itemize}

\textbf{AccDSL implementation contract (simplified):}

\begin{lstlisting}
event MeasureDefinedBenefit(
  plan: PensionPlan,
  dbo_opening: Money<C>,
  plan_assets_opening: Money<C>,
  current_service_cost: Money<C>,
  discount_rate: Rate<effective>,
  benefits_paid: Money<C>,
  actuarial_gain_loss_dbo: Money<C>, // From assumption changes
  actual_return_on_assets: Money<C>,
  contributions: Money<C>,
  period: YearFraction
) book EntityBook;

rule ias19_db for MeasureDefinedBenefit as e {
  // Service cost
  post {
    debit EmployeeBenefitExpense e.current_service_cost;
    credit DefinedBenefitObligation[e.plan] e.current_service_cost;
  }

  // Net interest
  let net_liability_opening = e.dbo_opening - e.plan_assets_opening;
  let net_interest = net_liability_opening * e.discount_rate * e.period;
  post {
    debit InterestCost_Pension net_interest;
    credit DefinedBenefitObligation[e.plan] (e.dbo_opening * e.discount_rate * e.period);
    debit PlanAssets[e.plan] (e.plan_assets_opening * e.discount_rate * e.period);
  }

  // Benefits paid
  post {
    debit DefinedBenefitObligation[e.plan] e.benefits_paid;
    credit PlanAssets[e.plan] e.benefits_paid;
  }

  // Contributions
  post {
    debit PlanAssets[e.plan] e.contributions;
    credit Cash e.contributions;
  }

  // Remeasurements to OCI
  let return_on_assets_excl_interest = e.actual_return_on_assets -
    (e.plan_assets_opening * e.discount_rate * e.period);
  post {
    debit PlanAssets[e.plan] return_on_assets_excl_interest;
    debit_or_credit OCI_Remeasurements_DB
      (e.actuarial_gain_loss_dbo - return_on_assets_excl_interest);
    credit DefinedBenefitObligation[e.plan] e.actuarial_gain_loss_dbo;
  }
}
\end{lstlisting}

\warn{Status: Not Implemented.} Simplified; missing:
\begin{itemize}
\item Actuarial assumption specification (mortality, turnover, salary growth)
\item Past service cost from plan amendments (19.103--19.108)
\item Asset ceiling (19.64)
\item Multi-employer and group plans (19.32--19.39)
\item Termination benefits (19.159--19.171)
\end{itemize}

\section{Income taxes: IAS 12}

\subsection{Deferred tax: temporary differences (IAS 12.5, 12.15--12.45)}

Deferred tax assets (DTA) and liabilities (DTL) arise from temporary differences
between carrying amounts and tax bases.

\textbf{Taxable temporary differences (12.15--12.31):} Future taxable amounts
$\Rightarrow$ DTL.

\textbf{Deductible temporary differences (12.24--12.33):} Future deductible amounts
$\Rightarrow$ DTA (if probable taxable profit).

\textbf{AccDSL implementation:}

\begin{lstlisting}
event RecognizeDeferredTax(
  asset_account: Account,
  carrying_amount: Money<C>,
  tax_base: Money<C>,
  tax_rate: Real in [0,1],
  probable_future_profit: Boolean
) book EntityBook;

rule ias12_deferred for RecognizeDeferredTax as e {
  let temporary_difference = e.carrying_amount - e.tax_base;
  let deferred_tax = temporary_difference * e.tax_rate;

  if temporary_difference > 0 { // Taxable temporary difference
    post {
      debit TaxExpense_Deferred deferred_tax;
      credit DeferredTaxLiability deferred_tax;
    }
  } else if temporary_difference < 0 && e.probable_future_profit {
    post {
      debit DeferredTaxAsset abs(deferred_tax);
      credit TaxExpense_Deferred abs(deferred_tax);
    }
  }
  // No recognition if DTA not probable
}
\end{lstlisting}

\warn{Status: Not Implemented.} Missing:
\begin{itemize}
\item Recognition exemptions (initial recognition, goodwill per 12.15, 12.21)
\item Tax loss carryforwards and credits (12.34--12.36)
\item Rate changes and their recognition (12.47--12.52)
\item Business combinations (12.19, 12.66--12.68)
\item Uncertain tax positions (IFRIC 23)
\end{itemize}

\section{Implementation roadmap and review protocol}

\subsection{Priority standards for P2/P3}

\begin{longtable}{L{35mm}L{55mm}L{46mm}}
\toprule Standard & Operations needed & Review milestone\\\midrule\endhead
IFRS 9 (instruments) & Full ECL staging, hedge accounting, embedded derivatives &
P3: Accounting firm review + differential testing vs. vendor systems\\
IFRS 15 (revenue) & Variable consideration, modifications, principal/agent,
licenses & P3: Worked examples across industries\\
IFRS 16 (leases) & Variable payments, modifications, sale-leaseback & P2: Common
operating and finance leases\\
IAS 19 (pensions) & Actuarial assumptions, plan amendments, multi-employer &
P4: Actuarial consultant + accounting review\\
IAS 12 (taxes) & Temporary differences, carryforwards, uncertain positions &
P3: Tax specialist review\\
IFRS 3 (combinations) & Purchase price allocation, goodwill, step acquisitions &
P3: M\&A accounting specialist\\
IAS 36 (impairment) & CGU identification, value-in-use, reversal & P3: Valuation
+ accounting review\\
IAS 21 (FX) & Functional currency, translation, net investment hedges & P2:
Multi-currency examples\\
IFRS 2 (share-based) & Grant-date FV, vesting, modifications & P3: Compensation
consultant\\\bottomrule
\end{longtable}

\subsection{Conformance test protocol}

For each operation marked \done{Reviewed}:
\begin{enumerate}
\item \textbf{Paragraph mapping:} Explicit standard paragraph(s) cited
\item \textbf{Worked example:} Known inputs, expected journals, balances, and disclosures
\item \textbf{Edge cases:} Boundary conditions, optional treatments, judgment areas
\item \textbf{Differential tests:} Compare with Big Four guidance or vendor implementations
\item \textbf{Independent review:} CPA or equivalent confirms compliance interpretation
\end{enumerate}

Before promotion from \warn{Pending Review} to \done{Reviewed}, all five steps
must complete. Operations marked \warn{Not Implemented} require design before testing.

\subsection{Gap acknowledgment and residual risk}

This chapter provides \textbf{design-level mapping} for a subset of IFRS. It does
not constitute:
\begin{itemize}
\item Certified accounting advice
\item A complete implementation of any standard
\item Authority to apply IFRS without additional judgment and review
\item Substitution for professional accounting support
\end{itemize}

Users relying on AccDSL for IFRS reporting must:
\begin{itemize}
\item Obtain independent accounting review for material transactions
\item Validate that scope and elections match their fact pattern
\item Document any departures or approximations
\item Retain ultimate responsibility for financial statement preparation
\end{itemize}
